\chapter{Security of critical IT infrastructures}
\label{chap:01}
\setlength{\parskip}{1em}

\section{Historical development of IT infrastructure}

The computing ecosystem and its service models have evolved gradually to on-premises and cloud paradigms, becoming increasingly commercialized to meet customer demands and scaling issues, as broadly represented in Figure \ref{fig:developmental_milestones_of_computing_environment}.

\begin{figure}[htbp]
    \centering
    \includegraphics[width=0.7\textwidth]{developmental_milestones_of_computing_environment.pdf}
    \caption{Developmental milestones of computing environments (VM - virtual machine; C - container; PaaS, IaaS, SaaS - external services): (1) Mainframe. (2) Cluster of servers. (3) Geographically distributed servers with virtualization (visually represented by color difference). (4) Hybrid approach with cloud services.}
    \label{fig:developmental_milestones_of_computing_environment}
\end{figure}

Initially, there were only mainframes that required dedicated teams for maintenance and performed at a highly mechanical level, so their bulkiness and slow speed enforced considerations for the adaptation of distributed computations \cite[Sec. 3.1]{lindsay2021evolution}. Fortunately, improvements in IT manufacturing significantly contributed to the emergence of dedicated servers that could be networked together in a modular way to provide scalability. Following this trend, virtualization moved beyond bare-metal deployments by allowing isolation layers to securely run multiple virtual machines (VMs) on shared hardware to significantly enhance the efficiency of resource utilization \cite{randal2021ideal}. Eventually, increasing demand for international connectivity and economic globalization naturally formed large-scale data centers with a network of geographically distributed servers in order to optimize content delivery and processing \cite[pp. 9–11]{lindsay2021evolution}.

Nowadays, organizations can offload many responsibilities and parts of their IT infrastructure to the cloud and service providers, which abstract physical machines and offer instant on-demand resource provisioning and management, thereby reducing operational burdens via "as-a-service" models with specialized offerings \cite[Sec. 3.3]{alghofaili2021secure}. Because of this, modern software engineers often operate within high-level abstractions where a deep understanding of underlying hardware-software interdependence is no longer a strict prerequisite for facilitating general-purpose tasks. Such a service-based outsourcing approach still requires integrating and coordinating operations across delegated areas (e.g., security policies, privacy compliance, and external monitoring). Moreover, certain strategic systems (i.e., governmental and military) and their data are prohibited from being hosted in the public cloud due to secrecy laws and regional restrictions \cite[Sec. 3.4]{alghofaili2021secure}.

\section{Approaches for complexity management}

As the sector of IT infrastructure has expanded, its complexity and associated challenges have emerged accordingly. Nonetheless, this complexity has become the primary factor for the process automation against manual, repetitive, and impractical operations. Therefore, automation and orchestration are no longer optional - they are critical for leveraging fully scalable and reliable systems from a strategic perspective. One of the natural strategies for complexity management is code. The process of codification captures procedures, routines, or algorithms of actions in a textual format \cite[pp. 1–2]{ouriques2023role}. For example, Chuprikov et al. propose specifying security policies in IT systems in the form of domain specific language (DSL) that can be run within an execution engine \cite{chuprikov2025security}. Codification also established the paradigm of Infrastructure-as-Code (IaC), which is focused on specifying and maintaining the state and content of IT infrastructures within self-documenting code (structured files with metadata, configurations, and procedures) with emphasis on provisioning reproducibility and transparent modifications \cite{pahl2025infrastructure}. This approach minimizes human errors and establishes immutable infrastructure using the phoenix model, allowing any server instance to be recreated from code without losing critical configurations. Thus, it eliminates the existence of snowflake servers, which are deviated systems with configuration drift and unpredictability due to their manual modification and a lack of documentation \cite[Ch. 2]{morris2020infrastructure}. 

IaC operations are managed by system administrators, DevOps teams (Development and Operations), SREs (Site Reliability Engineers), or IT professionals, as shown in Figure \ref{fig:iac_as_self_documenting_system}. They utilize IaC toolkits based on a spectrum of paradigms. Firstly, imperative one describes the process of achieving the target state via scripting languages and native shell scripts. Secondly, declarative one describes the final target state using definition files, configurations, and manifests. Usually, declarative ones (e.g., Terraform, Nix, Ansible) are implemented internally via general-purpose languages, but using configurations as application programming interfaces (APIs) \cite[Ch. 4]{morris2020infrastructure}. For example, containerization tools like Docker use Dockerfiles to define application dependencies and versions, while the runtime builds a container with artifacts according to specified requirements.

\begin{figure}[htbp]
    \centering
    \includegraphics[width=0.8\textwidth]{iac_as_self_documenting_system.pdf}
    \caption{Processes of applying changes to servers: (1) Manually or with simple scripts by following separate documentation. (2) IaC pipeline as a self-documenting system.}
    \label{fig:iac_as_self_documenting_system}
\end{figure}

Any IaC file should be proceeded through the continuous integration and continuous delivery (CI/CD) pipeline, which maintains safe deployment by runtime verification tests within simulated conditions for dynamic analysis and automated scans with validation checks of the codebase for static analysis \cite[pp. 9–10]{ecso2024technical}. According to Saavedra et al., incorporating intermediate representation (IR) and source-to-source translation into IaC static analysis facilitates the development of language-agnostic methods that can be ported to any IaC format \cite{saavedra2023polyglot}. So, IT specialists emphasize that IaC codebases should be maintained in accordance with secure coding standards on a regular basis (e.g., ISO/IEC series, NIST frameworks, etc.) \cite{konala2025framework}.

Sandobalin et al. examined the advisability of using IaC modeling tools by training control groups of computer science students on Argon and Ansible \cite{sandobalin2020effectiveness}. The research showed that Argon, as a visual toolkit, is easier to learn and can abstract complexity as a thin layer above existing IaC configuration formats \cite[pp. 14–15]{sandobalin2020effectiveness}. Thus, allowing users to simply draw their IT architecture, while the underlying system automatically generates a specific configuration for deployment. It implies that complexity management and computational thinking also extend to graphical and other representations of information. Diagrams and visualizations provide two-dimensional or multidimensional representations, enabling more structured modes of reasoning, whereas text is inherently a linear representation of information. Certainly, these forms should not be regarded as mutually exclusive, rather they synergize with each other and are interchangeable at different levels of abstraction.

In this regard, Model Driven Engineering (MDE) methodology proposes treating documentation, specifications, and model definitions as primary and reference artifacts describing behavior and components of a system, from which code generation toolkits can facilitate transformation between them (e.g., into source code templates in any programming language or pseudocode) in order to achieve synchronization \cite{verbruggen2023practitioners}. It bridges the gap between developers and business-oriented specialists by demonstrating certain domain-specific problems and rules, so they can focus more on business logic. For instance, Alfonso et al. developed the BESSER platform based on MDE, which provides a canvas editor and multiple design notations like UML (Unified Modeling Language) for generating backend applications from diagrams \cite{alfonso2024building}. Such toolkits usually support the development of reliable prototypes and concepts, but they are not intended for high-load or performance-critical systems. MDE also allows structures to dynamically adapt to emerging requirements instead of manually rewriting boilerplate code and thinking about platform-specific implementations, so the development process is optimized. For example, Chillón et al. implemented a CodeQL parser to automatically propagate component changes and enable migration between layers (database, logic, and services) by defining a single synchronizable schema instead of repetitive manual transformations, conversions, and modifications \cite{chillon2023propagating}. Additionally, AI in the form of generative language models enables rapid code generation from specifications as detailed prompts and instructions \cite{stoica2024specifications}. Although they can cause challenges due to hallucinations and their stochastic nature, it might be partially mitigated through enforcing strict rules, validation mechanisms, and multiple series of output refinements \cite{shankar2024who}.

\section{Resource management and accounting}

OWASP (Open Worldwide Application Security Project) identified 10 major risks for IT infrastructure, of which 5 are directly related to IT resources accounting \cite{owasp2025infrastructure}: "Outdated Software", "Insufficient Threat Detection", "Insecure Resource and User Management", "Insecure Access to Resources and Management Components", "Insufficient Asset Management and Documentation". This fact clearly emphasizes that proper management of all system components and organizational knowledge is mandatory. In this context, there is already a structured approach for collecting and accounting software components, such as software composition analysis (SCA), which defines common specifications and governance principles \cite{sabetta2024known}. It facilitated the establishment of a widely adopted format for software bill of materials (SBOM) with its multiple variations that allows to reduce a total vulnerability response time \cite{cisa2025shared}. Supporting this requires dedicated tools that collect data about software and infrastructure across various stages of their lifecycle. For example, researchers of \cite{klein2018infrastructure} developed an assisting toolkit that captures data from active production instances and further generates configurations or snapshots to recreate the environment on a different target machine or system. Usually, such procedures are combined with version control systems (VCS) that store and track versions of all project files, allowing teams to have process of collaborative development, a single source of truth, and well-defined compliance by localizing areas of responsibility \cite{ragavan2021version}. Specifically, SBOM and SCA provide an understanding of the connections and links between system components that form "dependencies", which directly and indirectly influence the functionality and quality of the base system, as demonstrated in Figure \ref{fig:components_of_abstract_complex_system}.

\begin{figure}[htbp]
    \centering
    
\includegraphics[width=0.7\textwidth]{components_of_abstract_complex_system.pdf}
    \caption{An abstract complex system (e.g., an operating system) consists of multiple underlying operational components and subsystems that provide its emergent qualities.}
    \label{fig:components_of_abstract_complex_system}
\end{figure}

According to statistics of the European Cyber Security Organisation (ECSO), solution-specific logic often represents only about 10\% or less of an application's total code, while the remainder is largely made up of third- party dependencies such as software (module, package, library), service (requests to external systems and network API), infrastructure (like drivers and execution platform) \cite[p. 8]{ecso2024technical}. Given a script in the Python programming language, the internal dependency in this case is the self-contained code, while the external dependency is everything that was created by third parties such as the implementation of standard library, interpreter/compiler, specification, OS, platform-specific requirements, and so on, as shown in Figure \ref{fig:python_ecosystem_from_perspective_of_dependencies}.

\begin{figure}[htbp]
    \centering
    \includegraphics[width=0.8\textwidth]{python_ecosystem_from_perspective_of_dependencies.pdf}
    \caption{Ecosystem of the Python programming language from the perspective of dependencies.}
    \label{fig:python_ecosystem_from_perspective_of_dependencies}
\end{figure}

From a broader perspective, the modern economy is service-oriented and highly interconnected in the global supply chain, where everyone relies on each other, as represented in Figure \ref{fig:simplified_example_of_global_supply_chain}. Hence, engineers can mitigate issues of resource constraints (e.g., schedule, personnel, risks, requirements, finance, and qualifications) by integrating external dependencies and reusing code as outsourcing for common problems during the implementation and maintenance of IT products.

\begin{figure}[htbp]
    \centering
    \includegraphics[width=0.8\textwidth]{simplified_example_of_global_supply_chain.pdf}
    \caption{Simplified example of the global supply chain in the context of IT.}
    \label{fig:simplified_example_of_global_supply_chain}
\end{figure}

\section{Supply chain attacks and dependency theory}

On the other hand, any risks and failures in a single link can cause a cascading effect across the entire chain, which is known as a supply chain attack. This attack vector exposes the ability to implicitly or explicitly influence on any IT infrastructure by leveraging direct and indirect (i.e., transitive) dependencies, invoking instability and fragility \cite[pp. 5–7]{hammi2023software}. It can be caused by any potential way of external influence, including unintentional flaws and defects, intentional threats like malicious software (malware), or destabilizing factors like breaking API changes and internal modifications \cite{okafor2022sok}. Inadequately governed dependency practices within a team can increase the risk of exposure to software supply chain attacks. Accordingly, it is necessary to be familiar with approaches to dependency management, such as:
\begin{itemize}
    \item Manual ("vendoring") - by copying an external dependency directly into the codebase in order to maintain and modify it independently \cite[p. 6]{cox2019surviving}.
    \item Automatic - dependency/package managers as a separate software system that is responsible for retrieving, handling, updating, and executing additional lifecycle operations on project dependencies \cite[p. 1]{cox2019surviving}.
\end{itemize}

These techniques are combined with build specification files (e.g., CMake, Make, Setuptools) and various distribution formats (e.g., wheels, source files, compiled binaries, and so on). Although some practitioners argue that the manual approach provides a more detailed understanding of dependency contents (source and distributions) and greater control (updating only explicitly), both of them are equally susceptible to supply chain attacks due to the propagation of transitive dependencies. Consequently, dependencies that are not explicitly pinned should not be relied upon, since they may disappear or be replaced within the direct dependencies from which they originate. Adversaries actively exploit such kind of characteristics and behaviors. Someone can abuse identity confusion by impersonating existing packages (by name similarity, description, or other metadata), claiming unregistered package names, and performing additional actions intended to mislead tools or users \cite{neupane2023beyond}. Therefore, it is essential to understand the underlying mechanisms, such as dependency resolution and metadata parsing. From a theoretical standpoint, dependency resolution/solving is a variation on the topic of the constraint satisfaction problem where the objective is to satisfy a predefined set of conditions and rules. Abate et al. reveal that it has proven to be an NP-hard problem for which no global optimal solution exists, so there are variety of different solvers and algorithms for this task \cite{abate2020dependency}. Ideally, such a solver should be decoupled as an independent module, regardless of the specific package ecosystem \cite[pp. 2–3]{abate2020dependency}. For example, the process of dependency management in the Python ecosystem relies on two components:
\begin{itemize}
    \item The "venv" module is used to store packages in a local virtual environment, preventing the pollution of the global environment and keeping dependencies separated per project.
    \item The "pip" package manager utilizes a backtracking dependency resolution algorithm to manage packages sourced from the PyPI (Python Package Index) repository, as shown in Figure \ref{fig:dependency_resolution_and_parsing}.
\end{itemize}

\begin{figure}[htbp]
    \centering
    \includegraphics[width=1.0\textwidth]{dependency_resolution_and_parsing.pdf}
    \caption{Process of dependency management: (1) Dependency resolution in Python-like pseudocode. (2) Dependency parsing.}
    \label{fig:dependency_resolution_and_parsing}
\end{figure}

In other ecosystems, NodeJS and its "npm" package manager address this process through workarounds by downloading multiple duplicate dependencies of varying versions, which increases the size of the "node\_modules" directory but prevents "dependency hell" where no valid resolution can be found \cite[p. 5]{abate2020dependency}, \cite{jafari2023dependency}. Generally, constructing the dependency tree or graph becomes increasingly complex with each additional dependency and constraint (version requirements or other ecosystem-specific rules), as illustrated in Figure \ref{fig:brief_introduction_to_dependency_theory}. Firstly, version constraints restrict the allowable range of package versions from the initial published version to the latest one according to the versioning format. Such versions basically represent a variation of the same entity. Secondly, dependency groups delimit or expand the number of required packages, because some components can be mandatory while others remain optional. Thirdly, Figure \ref{fig:brief_introduction_to_dependency_theory} represents only a basic example with a limited number of dependencies. In real projects, the number of dependencies and constraints is significantly larger, leading to dependency conflicts characterized by cross-references, circular or recursive relations, loops, duplicates, and similar pathological structures. These situations might be resolved, for example, by flattening the dependency tree into a more linear structure.

\begin{figure}[htbp]
    \centering
    \includegraphics[width=1.0\textwidth]{brief_introduction_to_dependency_theory.pdf}
    \caption{Elements of dependency theory: (1) Version constraints. (2) Dependency groups (concept from set theory). (3) and (4) Nested dependency graph (concept from graph theory) with direct and transitive ones (e.g., "1:M" as "one to many" relationships and other types of relations between them).}
    \label{fig:brief_introduction_to_dependency_theory}
\end{figure}

\section{Malicious threats and flaws in security}

Software engineers often prioritize functionality over security and make critical trade-offs under the assumption that security measures represent unnecessary performance overhead and justifying a lack of access restrictions to maintain a community that is “open to everyone” \cite{sasse2016debunking}. As a result, they potentially expose themselves and others to the risk that anyone may maliciously influence ecosystems. So, users are left to manage the external chaos on their own. Despite rising SCA and SBOM adoption by developers, researchers of \cite{alfadel2021empirical} indicate that supply chain incidents continue to increase, also noting that the remediation process in the Python ecosystem remains slow, with a median of 4 months required to fix all occurrences of a security vulnerability within the PyPI repository. Moreover, existing market solutions cannot be regarded as definitive protections against supply chain attacks, as experts identified several additional drawbacks and limitations of current SCA tools \cite{dietrich2023security,enck2022top,bohme2025software}:
\begin{itemize}
    \item focusing more on static analysis rather than on dynamic analysis during runtime execution;
    \item comprehensive deep scanning of transitive dependencies is often missing;
    \item insufficient integration with AI and data analytics-based systems.
\end{itemize}

Furthermore, obtaining reliable data about software remains challenging due to its various distribution formats and categories:
\begin{itemize}
    \item White box - fully known content (e.g., open source software).
    \item Gray box - partially known content (e.g., leaked information, rumors, reverse engineering, insider knowledge, alternative products, or forks).
    \item Black box - unknown content (e.g., functionality behind API like software-as-service, proprietary and closed-source software).
\end{itemize}

General SCA tools typically focus only on the white box approach by gathering already indexed vulnerabilities or components, thereby leaving critical gaps in threat visibility and transparency. Evidently, software analysis and detection/identification for other categories (both malware and legitimate) are inherently more complex and complicated due to their specifics, because anti-analysis techniques are used to protect such software and intellectual property, making component-level inspection difficult or reasonably impossible \cite{zaidenberg2022trulyprotect}. Although services like VirusTotal or similar antivirus platforms scan uploaded untrusted artifacts against numerous antivirus engines and databases to protect the corporate or personal sector, they cannot cover all possible forms of malware. For example, on platforms such as Google Play or similar repositories, there are many instances of hidden malware or fully undetectable malware (FUD) that have successfully bypassed the verification stage. Unfortunately, there is no universal solution or method against that, since software protection is constantly evolving. Therefore, it is mandatory to conduct additional reconnaissance and data gathering of existing samples and resources via well-known or standard deconstruction methods. In the realm of malware analysis, objects that are stored for threat intelligence are often termed as "IoCs" (Indicators of Compromise), "artifacts", or simply "samples" \cite{iklody2018decaying}. For example, these include entities such as IP addresses, URLs (links), DNS (domain names), virus signatures, hashes, files, and more. It is essential to clarify that the previously mentioned malware analysis is not limited solely to analyzing malware, and the term "malware" itself is quite vague/broad. Essentially, malware is any software which is considered potentially harmful or could lead to a security breach, depending on the specifics of our threat modeling with a predefined set of malicious characteristics and behavior. Malware analysis itself is typically categorized into four main types, which can be combined in various ways \cite[Ch. 3]{chenet2024survey}, \cite{ormeir2020dynamic}:
\begin{itemize}
    \item Static analysis focuses on examining the malware without executing it, which involves inspecting code, file structure, metadata, strings, and other non-executable aspects.
    \item Dynamic analysis, on the other hand, observes the malware in action by executing it in a controlled sandbox environment. This type of analysis monitors real-time / runtime behavior, such as network activity, file modifications, and registry changes, and is sometimes called "behavioral analysis".
    \item Manual analysis involves hands-on inspection of code or source by an analyst, allowing for in-depth, targeted evaluation.
    \item Automated analysis, on the other hand, leverages tools to systematically assess and classify malware or provide preliminary insights.
\end{itemize}

The fragmented software ecosystem worsens this situation, making accountability for the many packages and their adherence to standards unclear \cite{williams2025enabling}. For instance, Mirakhorli et al. reviewed 84 available SBOM management toolkits across multiple programming language ecosystems, which demonstrates the lack of unification and significant fragmentation in current solutions \cite{mirakhorli2024landscape}. Even a software package trusted by the community cannot offer proper guarantees, because mainstream software licenses clearly state that their products are distributed on an "as is" basis, meaning any issues that arise fall entirely outside their scope of responsibility \cite[pp. 2–3]{cox2019surviving}. Furthermore, software licenses contain other ambiguities \cite{wintersgill2024law}:
\begin{itemize}
    \item They have limited legal power and lack formal enforcement in certain countries, with cases of violations.
    \item Transitive dependencies may have conflicting licenses compared to the primary dependency, requiring the use of scanning tools to account for overlapping permissive and restrictive licenses.
\end{itemize}

Experts propose different solutions for the current status quo, ranging from weak to radical measures. Firstly, there should be an introduction of increased mandatory and strict basic requirements ("baselines") at the platform level with compliance from end user side. This could reduce the number of attacks and malicious code due to high standards and improved quality. It also reflects the principle of "secure by default", where secure policies and good default configurations should become common for many, and not just the exception reserved for professionals:
\begin{itemize}
    \item Mandatory oversight and analysis before package publication and subsequently with every update/change. For example, the establishment of a moderation system where the community will naturally report and react to the appearance of malicious packages. Such movements as bug bounty programs can provide payment to the community for finding bugs to establish natural vulnerability research and embrace ethical hacking over illegal activities \cite{bozzini2025regimes}. Also, there is a productive case within Astana IT University where they regularly implement CTF (capture the flag) and cyberpolygons to train security engineers, thereby validating skills and preparing professional specialists for organizations \cite{abdiraman2023comparative}.
    \item Policies for temporary blocking or issuing "call down" periods between user actions during anomalous activity (e.g., IP address change, password or specific token renewal, recent publications, and so on).
    \item Movement towards slowing down of overproduction and starting to "recycle waste" by rethinking legacy systems or abandoned ideas (re-evaluating them, finding patterns, and recognizing repetition) \cite{assuncao2025contemporary}. Possibly, developers should also strive to regularly improve the standard library and make it better with "batteries included" in the official distribution. Otherwise, its position will be taken by many competing dependencies covering the exact same specific topic, thereby generating a high degree of entropy.
    \item Designing systems in a way that they are "secure by design" in their foundation means that security and safety are built in initially, rather than being an additional element. This approach dictates that engineers strive to place security and features at least on the same level of value during the development process \cite{pohl2015secure}. It eliminates specific categories of vulnerabilities, restricts/limits malware, and forms a multilayered security / defense-in-depth approach \cite{mell2016measuring}.
    \item Next-generation and future-proof systems that are comprehensive to the extent that they can self-analyze and self-improve, making them more reliable and less fragile. For instance, Berhe et al. showed that companies are compelled to have a mitigation plan and a regular budget for automated software dependency recovery or manual actions such as replacing it with alternatives, upgrading it, or reverting to a previous release \cite{berhe2023maintenance}. Ideally, these issues could be mitigated by well-decomposed systems that might continue to function in such critical cases or have internal recovery mechanisms ("self-healing") for the failure of individual components only up to a certain limit \cite{dehraj2021review}.
\end{itemize}

Secondly, the software market should pay more attention and show greater concern for formal regulatory mechanisms to increase trust and responsibility in the community, such as certification, auditing, and transparent processes, along with initiatives aimed at improving standards. Sammak et al. attended an interview among a small subset of experienced developers, which revealed that the majority of respondents are unfamiliar with details of software supply chain management or have a superficial understanding due to the absence of proper training on this topic for a broader audience \cite{sammak2023developers}. Thirdly, popular or operation-critical software products have a strong influence on the whole landscape, so compromising them can undermine not just dependent infrastructure but also the global contracts of security, such as the cryptography stack \cite{tsoupidi2025protecting}. For example, self-replicating malware or a virus might compromise a significant portion of the IT ecosystem by ensuring that any code compiled on an infected hardware platform, compiler, or OS becomes compromised as well \cite{ladisa2022taxonomy}. This is due to the fact that definition files are merely instructions without guarantees that our intentions will be correctly expressed and transformed with preserved semantics in the produced output (e.g., machine code, process, and other representations) by an intermediary system (compiler, interpreter, execution environment) \cite{gottschlich2018three}, as shown in Figure \ref{fig:intermediary_system_between_design_and_implementation}. Furthermore, conducting a compromise assessment or software evaluation for the produced product is complicated, because infected reverse engineering tools (e.g., debugger, decompiler, disassembler) or even their clean versions running in a compromised environment might hide the infected parts and backdoors as blind spots in the software in order to remain undetected.

\begin{figure}[htbp]
    \centering
    \includegraphics[width=1.0\textwidth]{intermediary_system_between_design_and_implementation.pdf}
    \caption{An intermediary system between definition and implementation.}
    \label{fig:intermediary_system_between_design_and_implementation}
\end{figure}

\section{Trust factors and social aspects}
Due to the high stakes, it is crucial to determine whom to trust and based on what principles to minimize such risks beforehand. Researchers Singer and Bishop propose viewing trust in the form of a relationship model that makes us vulnerable because we rely and depend on a third party without any doubt due to assumptions and distractions, considering them trustworthy \cite{singer2020trust}. It means that our inherent tolerance to possible flaws/mistakes due to the factor of trust affects on our cybersecurity posture. They also suggest applying threat modeling and risk analysis before and during the design stage in the development lifecycle to minimize the factor of trust \cite[Ch. 4]{singer2020trust}. Hence, trust should not be considered as a static and accurate indicator, but rather it is a dynamic probability of being trusted by someone else and vice versa, where any mistake from either party leads to its loss or reduction. For example, one of the parties (recipient or sender) can at any time disrupt data transmission and its correctness by intentionally beginning not to adhere to the given standard protocols and interaction interfaces.

For these reasons, the National Institute of Standards and Technology (NIST) developed the concept of Zero Trust Architecture (ZTA), in which trust requires continuous verification \cite{rose2020zero}. Babenko et al. clearly demonstrated that empirical data on ZTA implementation correlates with a lower incident rate of insider threats \cite{babenko2024cybersecurity}. Organizations can adopt this by following the key principles listed below:
\begin{itemize}
    \item Continuous evaluation of all actions and actors instead of granting full access within a protected perimeter. This means the environment itself does not affect how we trust someone/something, assuming operating in a hostile environment \cite[Ch. 2.2]{rose2020zero}.
    \item Eliminating implicit trust in favor of explicit trust using a trust algorithm as a method for verifying access or credibility with transactional properties \cite[Ch. 3.3]{rose2020zero}.
    \item Segmentation of the network and systems into atomic cells, which theoretically fragments the attack surface (e.g., a large monolith will expose more data than individual microservices), as visualized in Figure \ref{fig:network_segmentation_in_zero_trust} \cite{cisa2025journey}.
    \item Least privileges of access and minimizing the amount of dependencies via strict comprehensive assessments of vendors or suppliers \cite{collier2021zero}. On top of this, it is recommended to apply DevSecOps practices for the automation of software building, testing, and the overall "shift-left" movement for the early identification of problems in information security \cite{dodcio2024enterprise}.
    \item Security as a dynamic process and adaptive system is achieved through the continuous collection and analysis of data/statistics (resources, assets, incidents, etc.). In particular, Mushtaq et al. note the importance of compliance with GDPR (General Data Protection Regulation), explainable AI, and anomaly detection to maintain transparency in ZTA \cite{mushtaq2025systematic}.
\end{itemize}

\begin{figure}[htbp]
    \centering
    
\includegraphics[width=1.0\textwidth]{network_segmentation_in_zero_trust.pdf}
    \caption{Network segmentation in ZTA, where each segment consists of multiple defense layers and granular microservices.}
    \label{fig:network_segmentation_in_zero_trust}
\end{figure}

One example of ZTA implementation is decentralized systems (e.g., blockchain), which contain zero-knowledge proofs and consensus algorithms such as "proof of work" and others that guarantee trust based on a defined trust factor (e.g., computational power, number of devices/actors, amount of specific tokens) \cite{bangui2024when}. Furthermore, there is a concept of "trustless" in the same sphere that further reinforces ZTA and trust management approaches in distributed systems. Certain studies even consider applying properties of distributed systems and blockchain to version control systems and secure package management, leveraging their data integrity and non-repudiation characteristics to enhance supply chain security \cite[pp. 8–12]{hammi2023software}, \cite{grilli2024combining}.

Overall, the current landscape of security for critical IT infrastructure is broad and continues to evolve alongside the increasing sophistication of cyberattacks. It is now evident that there is a significant need for more robust dynamic analysis of software and its underlying components, as static source code analysis and standard tooling are insufficient. Future large-scale or critical incidents will likely serve as catalysts, compelling the global community to fundamentally rethink and restructure prevention strategies. So, the industry is going to shift toward proactive prevention, the establishment of comprehensive security protocols, and the rigorous verification and analysis of IS.

\section{Distributed computing and blockchain-based systems}

The development of the Internet and the advancement of networking technologies provided capabilities for computing systems to operate across geographically distributed locations, accelerating information exchange, the accumulation of collective knowledge, and hosting both public or private digital services \cite[pp. 9–10]{leiner2009brief}. For example, crypto networks as overlay networks and Web3 technologies allow to conduct direct financial transactions and peer-to-peer (P2P) connections between multiple nodes. Several of these decentralized networks leverage blockchain as an immutable data structure to ensure trust and data integrity \cite[Ch. 3]{hassija2021framework}. Besides, there are alternative architectures and structures (e.g., graphs, trees, unstructured ones, and so on \cite[Ch. 6]{hassija2021framework}) with other characteristics and without known limitations of blockchain:
\begin{itemize}
    \item Versioning or migration during hard forks requires branching the chain. It remains a voluntary decision for network participants, as they may choose to stay on the legacy branch or migrate to the new one \cite[pp. 1–2]{conlon2024bitcoin}.
    \item Scaling of nodes requires off-chain data separation, as large blocks cause network congestion and degrades transaction performance \cite[Ch. 2–3]{chen2024comprehensive}.
    \item Block propagation or relaying time requires prevention of redundant data transmissions and shorter distances to be constant, which is difficult to maintain across large distances between nodes when also dealing with cross-border deployments or regional restrictions \cite{zhang2022speeding}.
\end{itemize}

Distributed systems are characterized by having consensus protocols and zero-knowledge proofs, which are both mechanisms for validating or verifying computational correctness and compliance with a predefined set of rules \cite{lavin2024survey}. Their development and approbation are complex processeses involving rigorous vulnerability analysis and formal verifications \cite{tang2024zeroknowledge}. For example, the framework proposed in \cite{jones2026towards} integrates proof assistants with generative language models to automate the mathematical verification of protocols. However, experts raise ecological concerns by shifting towards cheap renewable energy and energy-efficient optimizations, because certain consensus mechanisms such as Proof of Work (PoW) require intensive redundant computations that consume more power for protocol security than providing tangible utility for humanity \cite{alzoubi2023green,tomatsu2023bitcoin}. Therefore, \cite{hoffmann2022challenges} and \cite{lihu2020proof} advocate for Proof of Useful Work (PoUW), which provides practical applications for AI model training, algorithmic optimization, data labeling, and other valuable tasks. Moreover, various factors within certain crypto networks (e.g., dynamic difficulty, halving events, and market hype) have shaped their economics to token holding and the centralization of power within large pools. For instance, FruitChains tries to mitigate it by guaranteeing fair reward distribution proportional to wasted computational power \cite{pass2017fruitchains}. Also, Auditable Proof-of-Work (APoW) allows nodes to cryptographically prove their efforts, audit participants, and detect distributed attacks \cite{lerner2026apow}.

Centrally controlled and widely used cloud services (e.g., Amazon Web Services, Google Cloud, Microsoft Azure) can pose a significant risk to dependent infrastructure and users because their outages and failures can halt essential operations, as shown in the aggregated data of Talluri et al. from 2018 to 2020 \cite{talluri2026cloud}. The authors also state that these situations are frequent and natural occurrences ranging from minor to major ones, noting that a discrepancy between official reports and often-unreliable user feedback makes it more difficult to measure the full extent of such incidents \cite[pp. 5–7]{talluri2026cloud}. Distributed ledger mechanisms can decentralize traditional IT infrastructures to improve reliability and uptime while reducing the risk of single points of failure and dependence on centralized management. However, Balduf et al. daily crawled traffic and network participants in the IPFS (InterPlanetary File System), one of the largest decentralized P2P storage, during a single month period in 2023, showing that 20300 (79.6\% of nodes) of approximately 30000 nodes across 152 countries were hosted on cloud-based VPS (Virtual Private Servers), with a notable concentration in the USA (33\% of nodes) \cite{balduf2023cloud}. It reveals that despite having decentralized networks, they still remain deeply interconnected with the underlying centralized ecosystem, while there are clear and understandable reasons for this phenomenon:
\begin{itemize}
    \item Difficulties of exposing a local node with public IP address due to the lack of IPv6 support and NAT (Network Address Translation) configuration from the ISP (Internet Service Provider) \cite[p. 6]{balduf2023cloud}
    \item The limited number of IPFS gateways acts as bottleneck, with their resolution partially reliant on Cloudflare services \cite[pp. 3, 11]{balduf2023cloud}.
    \item Usage of cloud services for privacy and convenience.
\end{itemize}

\section{Optimizing computing and small models on edge devices}

While historical trends indicate that computational demands have continuously increased since the inception of deep learning and throughout the current large-scale era of AI \cite[Ch. 3]{sevilla2022compute}, the continuous scaling paradigm is hitting a plateau. Physical constraints and diminishing returns make this approach increasingly unsustainable, requiring search for more extensive methods \cite{varoquaux2025hype}. Experts propose various modern solutions and enhancements from both hardware and software sides. For instance, Li et al. demonstrated that offloading partial neural network inference or other computational tasks from data centers to edge devices can reduce environmental impacts (e.g., energy consumption, carbon emissions, and water usage) in their experiments \cite{li2025case}. Particularly, it is possible to construct a pipeline/workflow that evaluates the complexity of a given task and redirects it to a smaller model, achieving reasonable generation speed and sufficient output quality for specialized tasks or domains \cite{belcak2025small}: mathematical proofs, seq2seq (sequence to sequence), NER (named entity recognition for structured extraction of entities or actions), classifications, formatting and autocompletion of coding snippets, automating repetitive actions with agentic AI, searching and sorting unstructured data, and so on. These capabilities and efficiency of on-device processing are primarily caused by the SoC (System-on-Chip) architecture of mobile devices, which allows utilization of GPU (GraphicsProcessing Unit), CPU (Central Processing Unit), and NPU (Neural Processing Unit) for tensor and matrix computations on a single chip \cite{lee2019ondevice}. Moreover, Amatuni et al. represent USA statistics in which the average lifetime/lifecycle of such devices can be extended by around 18\% (from 2.7 to 3.2 years) through secondary markets, thereby reducing the production volume of new units \cite{amatuni2026secondhand}. However, the root cause of the problem arises from manufacturers and distributors who refuse to make their devices repairable, reliable, durable, and equipped with interchangeable components across product lines or generations for ease of maintenance and repairment. Researchers \cite{oellermann2022open} and \cite{kalapothas2023survey} discuss that open hardware and electronics can democratize high-performance computing, enhance reproducibility, and facilitate the sharing of experience for iterative refinement by applying RISC-V instruction set, reprogrammable FPGAs (Field-programmable Gate Array), Vulkan and OpenGL backend instead of proprietary CUDA, ASIC (Application-Specific Integrated Circuit) as special-purpose hardware, and customizable microcontrollers. From the software part, there should be a reduction of power consumption and environmental footprint that involves optimizing algorithms (following Big-O notation), adopting energy-efficient programming languages (compiled languages instead of interpreted ones), increasing hardware utilization (concurrency and parallelism), lower memory usage of a programming language (PL), and restarting programs on demand instead of background running \cite{pereira2017energy}. Kempen et al. noted that specific implementations of PL with the same syntax can differ in their runtime, clearly showing that LuaJIT and PyPy outperform Lua and CPython by 500\% and 125\% in execution speed and energy efficiency respectively in their benchmarks, because JIT (Just-in-time) optimizer/compiler can simplify hotspots, loops, and frequent instructions \cite[pp. 7–8]{vankempen2025not}.

DeepSeek, an independent Chinese AI laboratory, demonstrated that training a 671B (billion) parameter state-of-the-art (SOTA) large language model can be achieved with 2.788 million hours of NVIDIA H800 GPU (approximately 5.6 million USD at 2024 market rates) through the use of knowledge distillation techniques \cite{deepseekai2025deepseekv3}. However, such estimates do not entirely represent the actual situation, as they do not account for associated costs such as:
\begin{itemize}
    \item maintenance of IT infrastructure and hardware accounting (e.g., service lifetime, wear, and degradation) due to overheating or throttling under high-intensity loads or environmental hazards;
    \item employee salaries and fees for external services;
    \item expanding density and scale of digital information for multimodal capabilities (video, audio, images, etc.).
\end{itemize}

In this regard, federated learning allows to shift the responsibility for physical hardware maintenance from a centralized server to a distributed network of nodes that serve abstraction of resources for AI-related tasks such as training models, processing large datasets, running model inference, data annotations, and so on \cite{yurdem2024federated}. It is categorized by several sub-types on data organization and structures:
\begin{itemize}
    \item Horizontal (same datasets but different samples), vertical (complementary features from overlapping samples), transfer learning (adapting a model trained on one dataset to perform a different but related task) \cite[Ch. 2–3]{yurdem2024federated}.
    \item Local (on a single device), decentralized (each node acts independently), personalized (selecting specific nodes based on their state or criteria), and centralized (as demonstrated in Figure \ref{fig:federated_and_distributed_learning}) \cite[Sec. 2.1-2.3]{li2020federated}.
\end{itemize}

\begin{figure}[htbp]
    \centering
    \includegraphics[width=1.0\textwidth]{federated_and_distributed_learning.pdf}
    \caption{Federated learning and distributed computations: (1). Centralized form of federated and distributed learning that are constrained by the networking limitations of endpoints on data volume and communication speed. (2). A worker node operates on heterogeneous hardware such as microcontrollers, single-board computers, and other devices.}
    \label{fig:federated_and_distributed_learning}
\end{figure}

\section{Security and privacy of execution environments}

Federated learning can also handle privacy aspects by having each participant share only a snapshot of their locally trained model (its insights in model weights and intermediate activation vectors) to aggregate and combine them via secure multi-party computation, rather than transferring raw datasets directly \cite{xu2023federated}. In the context of medical records, it is mandatory to implement data anonymization, model encryption, and differential privacy through noise injection \cite{koutsoubis2025privacy}. Also, it is even possible to facilitate ephemeral data access, in which data is temporarily cached in memory for processing and then immediately cleared \cite[p. 11]{xu2023federated}. Moreover, certain networks like Monero are strictly focused on participant privacy, private payments, private computations, and function privacy \cite[pp. 7–8]{almashaqbeh2022sok}.

Security in distributed systems is critical for stabilizing nodes and preventing localized issues from propagating across the network. Daulay et al. investigated the security risks of inference attacks aimed at extracting sensitive data from model weights and disrupting network integrity (e.g., data spamming), noting that the inherent complexity of distributed governance makes ensuring honest participation and data protection particularly challenging \cite{daulay2025novel}. Therefore, there is the concept of Trusted Execution Environments (TEEs) that introduces:
\begin{itemize}
    \item Hardware and software isolated execution restricts operations to an isolated environment, functioning as both a trusted and controlled environment \cite{jauernig2020trusted}.
    \item Mitigation of memory corruption, various vulnerabilities, and low-level kernel exploits by utilizing memory safety and specific modifications \cite{wan2020rustee}.
    \item Confidential computing for sensitive operations ensures that network participants remain auditable, protected, and verifiable \cite{ye2024position}.
    \item Application of TPM (Trusted Platform Module) chips for hardware-based trust, MAC (Medium Access Control) addresses for networking hardware identification, and specialized firmware to ensure the integrity of nodes \cite{bartock2022hardware}.
\end{itemize}

By default, mobile operating systems such as Android already offer robust and secure-by-default capabilities through restrictive application sandboxing and a strict permission model \cite{tuncay2024android}. Nevertheless, malware detection and dependency analysis still remain essential for evaluating untrusted applications and third-party libraries \cite{amro2017malware}. Chen et al. evaluated WebAssembly (WASM) as a viable execution runtime that provides a secure and platform-agnostic sandbox \cite{has2025webassembly}. Also, experts suggest using WASM for smart contracts, where programs are compiled or encoded into WASM bytecode \cite[Ch. 3]{rodriguezsela2025towards}. Although bad actors may exploit distributed networks for malicious purposes (e.g., Command and Control (C2) infrastructure, illegal operations, or bypassing protective measures to abuse resources), researchers propose methods for preserving the computational correctness and moderation of smart contracts by using fuzzers or scanners to identify vulnerable or malicious patterns \cite{chen2022wasai}.
